%%=============================================================================
%% Conclusie
%%=============================================================================

\chapter{Conclusie}%
\label{ch:conclusie}

% TODO: Trek een duidelijke conclusie, in de vorm van een antwoord op de
% onderzoeksvra(a)g(en). Wat was jouw bijdrage aan het onderzoeksdomein en
% hoe biedt dit meerwaarde aan het vakgebied/doelgroep? 
% Reflecteer kritisch over het resultaat. In Engelse teksten wordt deze sectie
% ``Discussion'' genoemd. Had je deze uitkomst verwacht? Zijn er zaken die nog
% niet duidelijk zijn?
% Heeft het onderzoek geleid tot nieuwe vragen die uitnodigen tot verder 
%onderzoek?

Deze bachelorproef onderzocht de risico's van log-gebaseerde replicatie in Db2 voor z/OS binnen kritieke financiële omgevingen met mainframes.
Hierbij ging er bijzondere aandacht naar de impact op Recovery Point Objective (RPO) en Recovery Time Objective (RTO) in geodistribueerde architecturen.
Uit de analyse blijkt dat een opstelling met uitsluitend twee synchrone datacenters, hoewel deze een RPO van nul garandeert, onvoldoende bescherming biedt tegen geregionaliseerde uitval. 
De afstandsbeperkingen inherent aan synchrone replicatie maken het onmogelijk om volledige geografische spreiding te realiseren zonder significante impact op de prestaties.
Asynchrone replicatietechnologieën, zoals GDPS Global Mirror, bieden daarentegen wel de mogelijkheid tot geografische spreiding over grote afstanden, maar introduceren onvermijdelijk risico's op dataverlies en inconsistentie. 
Deze risico's manifesteren zich voornamelijk onder de vorm van replicatievertraging, incomplete Logical Units of Work (LUW's) en inconsistenties in de volgorde van transacties. 
Deze fenomenen hebben een directe negatieve impact op zowel RPO als RTO en vereisen bijkomende herstelmechanismen.
De hybride architectuur, waarbij twee synchrone sites worden gecombineerd met een derde asynchrone site, vormt op basis van dit onderzoek de meest robuuste oplossing. 
Deze opstelling combineert de voordelen van beide replicatiemodellen: een RPO van nul en snelle failover binnen de synchrone regio, aangevuld met bescherming tegen regionale uitval via een geografisch gescheiden derde site. 
De keerzijde is echter een verhoogde complexiteit, vooral bij failover naar de asynchrone site, waarbij consistentieherstel en replicatieachterstand een significante invloed hebben op de RTO.
Binnen logreplicatie blijkt Controlled Log Catch-Up (CLC) een essentieel mechanisme om transactionele consistentie te waarborgen na falen. 
Hoewel CLC effectief is in het herstellen van consistente toestanden, verhoogt het de hersteltijd wanneer grote achterstanden moeten worden ingehaald. 
Hierdoor ontstaat een duidelijk spanningsveld tussen dataconsistentie en herstelsnelheid.
De bijdrage van dit onderzoek ligt in het systematisch in kaart brengen van de interactie tussen replicatiestrategieën, faalscenario's en hun impact op RPO en RTO binnen Db2 voor z/OS-omgevingen.
Door deze analyse worden financiële instellingen ondersteund bij het maken van onderbouwde architecturale keuzes die voldoen aan strikte regelgeving zoals DORA vanuit de Europese Unie.
Desondanks kent dit onderzoek beperkingen. 
De analyse is hoofdzakelijk gebaseerd op literatuur en theoretische modellering, zonder empirische validatie in een productieomgeving. 
Factoren zoals specifieke hardwareconfiguraties, netwerkcondities en workloadvariaties kunnen in de praktijk een andere impact hebben op de prestaties en risico's.
Verder onderzoek kan zich richten op experimentele validatie van de beschreven architecturen, bijvoorbeeld via simulaties of testomgevingen. 
Daarnaast is er ruimte voor onderzoek naar geautomatiseerde failovermechanismen en optimalisatietechnieken die de impact van replicatievertraging en CLC op RTO verder kunnen beperken. 
Ook de rol van nieuwe technologieën, zoals AI-gedreven monitoring en voorspellende faal analyse, vormt een interessant perspectief binnen dit domein.
Samenvattend kan gesteld worden dat er geen universele oplossing bestaat die simultaan een RPO van nul, minimale RTO en volledige geografische spreiding garandeert. 
De keuze voor een replicatiestrategie blijft steeds een afweging tussen consistentie, performantie en operationele complexiteit. 
De hybride architectuur biedt hierbij de beste balans voor kritieke systemen binnen financiële omgevingen, mits de bijkomende risico's actief worden geminimaliseerd door de besproken mitigatiestrategieën.